%!TEX program = lualatex
\documentclass[ngerman,headheight=70pt]{scrartcl}
\usepackage[ngerman]{babel}
\usepackage{blindtext}
\usepackage[utf8]{luainputenc}
\usepackage{scrlayer-scrpage}
\usepackage[T1]{fontenc}
\usepackage{fontspec}
\usepackage{multicol}
\usepackage{enumitem}
\usepackage{epigraph}
\usepackage[german=quotes]{csquotes}
\usepackage[bottom=2cm,footskip=8mm]{geometry}
\setmainfont{Source Sans Pro}
\MakeOuterQuote{"}

\setlength{\parindent}{0pt}
\setlength{\parskip}{6pt}

\pagestyle{plain.scrheadings}
\thispagestyle{scrheadings}
\lohead{\vspace*{-1.5cm}\includegraphics[scale=0.4]{../images/stupa_logo_rot.pdf}}
\rohead{\\Präsidium des\\Studierendenparlaments}

\lofoot[Universität Hamburg · Präsidium des Studierendenparlaments\\\hrule Von-Melle-Park 5 · D-20146 Hamburg]{Universität Hamburg · Präsidium des Studierendenparlaments\\\hrule Von-Melle-Park 5 · D-20146 Hamburg}
\cofoot[]{}
\cohead[\thepage]{}

\renewcommand{\thesubsection}{TOP \arabic{subsection}}
\renewcommand{\thesubsubsection}{\arabic{subsubsection}.}
\setcounter{subsection}{-1}

\begin{document}
    UHH · StuPa-Präsidium · Von-Melle-Park 5 · D-20146 Hamburg

    \section*{Beschluss des Studierendenparlamentes vom 1. Dezember 2016}
    \subsection*{Nach wie vor: Aufklärung, Kultur, Solidarität!}
    \subsubsection*{Aktionstag gegen Rechts und Veranstaltung am 30. Januar 2017}

    \blockquote[Thomas Fricke: "Was Europa von Amerikas Großen lernen kann."
    Spiegel Online, 11.11.2016]{
    Womöglich bräuchte Europa besser heute als morgen einen so eindrucksvollen
    New Deal, wie ihn Roosevelt damals gegen Krisenfolgen, Wutbürger und
    Absturzängste einsetzte.}

    \blockquote[Thomas Mann, „Spanien“, 1937]{
        Nun, das Geistige, unter dem politisch-sozialen Gesichtswinkel gesehen,
        ist das Begehren der Völker nach besseren, gerechteren, glücklicheren,
        dem Stande des menschlichen Selbstbewußtseins richtiger angepaßten
        Lebensbedingungen, -- es ist dieses Begehren in seinem Bejahtsein durch
        alle, die guten Willens sind.
    }

    Immer größer wird der Widerspruch zwischen Arm und Reich und die Konflikte
    um die gesellschaftliche Entwicklung spitzen sich weltweit zu. Der immense
    gesellschaftliche Reichtum konzentriert sich in den Händen Weniger, welche
    sich im Dauerkampf untereinander und gegen die Bevölkerung befinden. Millionen
    Menschen sind auf der Flucht vor Krieg, Armut und Umweltzerstörung und auch
    in den reichen Ländern wächst die Zahl der prekär Beschäftigten und
    Erwerbslosen. Vor diesem Hintergrund sind die Herrschenden und die politischen
    Opportunisten, die diese Entwicklung willfährig betrieben haben, in tiefer
    Legitimationskrise. Der Neoliberalismus ist wie Saturn, er frisst seine eigenen Kinder.

    Dass Donald Trump nächster Präsident der USA wird, ist vor allem Ausdruck
    dieser Umbruchsituation. Er steht mit anderen wie Le Pen, Wilders und Petry
    für jene Rechtspopulisten und Reaktionäre, die mit einer Anti-Establishment-Rhetorik
    Rettung versprechen, um mit ihrer misanthropen Ungleichheitsideologie und -praxis
    die Übel weiter zuzuspitzen.

    Ihnen steht alternativ eine wachsende Bewegung für Frieden, internationale
    Solidarität, soziale Gerechtigkeit, Erweiterung der demokratischen
    Partizipation, Umweltschutz und für ein aufgeklärtes menschliches Miteinander
    gegenüber. (Der mit dieser Bewegung verbundene Bernie Sanders hätte allen
    Umfrage zu Folge die Wahl gegen Trump mit Abstand gewonnen).

    In dieser globalen politischen Polarisierung ist eine progressive politische
    Entwicklung in Europa und der Bundesrepublik von enormer zivilisierender
    Bedeutung. Vor diesem Hintergrund gilt erweitert, was das Studierendenparlament
    mit seinem Beschluss vom 16. April 2016 formulierte: Die Verfasste
    Studierendenschaft der Uni Hamburg "sieht sich deshalb in der Pflicht, ihre
    gesellschaftliche Aufklärungsarbeit gegen Rechts, ihr Engagement für Frieden,
    Abrüstung, sozialen Fortschritt und für umfassende Demokratie zusammenzuführen
    und zu intensivieren. Das Studierendenparlament will diese Orientierung forcieren."

    Das Studierendenparlament greift daher seinen damaligen Beschluss auf und
    ruft einen Aktionstag gegen Rechts zu Beginn des Sommersemesters 2017 aus.
    Dieser soll sich besonders der analytischen, geschichtsbewussten und
    kulturkritischen Auseinandersetzung mit der AfD und anderen rechtspopulistischen
    Akteuren widmen und über eine wirkliche, soziale und (humorvoll) aufgeklärte
    Alternative zum neoliberalen Gesellschaftsmodell informieren und mobilisieren.

    Um auch im Wintersemester 16/17 noch zur Aufklärung und Mobilisierung gegen
    Rechts beizutragen, strebt das StuPa für den 30. Januar 2017 eine zusätzliche
    Veranstaltung an. An diesem Tag jährt sich zum 84. Mal die Machtübertragung
    an die deutschen Faschisten.

    Nur etwas mehr als einen Monat später trat in den USA der damals neu gewählte
    Präsident Franklin D. Roosevelt sein Amt an. Er repräsentiert, was damals
    positive Entwicklungsmöglichkeit noch innerhalb des Kapitalismus aus der Krise
    der 20er und 30er Jahre heraus gewesen ist. Mit massiven Investitionen in
    Infrastruktur und Arbeit, mit Mindestlöhnen, Rentenversicherung und
    Sozialgesetzen sowie Regulierung der Banken wurde in den Dreißiger Jahren in
    den USA ein Pro¬gramm zur Verbesserung der Lebensbedingungen durchgesetzt,
    das auch heute orientierend sein kann.

    Aus der Geschichte lässt sich lernen.

    Das Studierendenparlament beauftragt seinen Ausschuss gegen Rechts damit, die
    Veranstaltung und den Aktionstag gegen Rechts vorzubereiten und zu koordinieren.

    Für die Gestaltung des Aktionstages regt das Studierendenparlament Folgendes an:

    \begin{enumerate}
        \item Eröffnung mit einer zentralen Diskussionsveranstaltung zu aktuellen
        Erkenntnissen aus der historischen Faschismuserfahrung. Es möge eine Person
        gewonnen werden, die zum Einstieg referiert zur Verwandtschaft zwischen
        historischer vorfaschistischer Krise und heutiger Krise, vor allem aber
        zu Unterschieden und Lehren, die eine progressive Wendung der heutigen
         Krise möglich und wahrscheinlich machen.

        \item Hochschulgruppen und andere Aktive werden aufgerufen, über den Tag
        verschiedene Veranstaltungen und Workshops durchzuführen, die in einem
        gemeinsamen Programmheft zusammengeführt werden. Hierzu können Diskussionen
        und Vorträge gehören (z.B. zu Aktualität des Potsdamer Abkommens, zu
        1968 –- gegen den Muff aus 1000 Jahren unter den Talaren, zu internationalen
        Erfahrungen antifaschistischen Widerstands, zur aktuellen Auseinandersetzung
        mit Rechtspopulisten, Nationalisten und Chauvinisten in anderen Ländern,
        zur Identitären Bewegung und Burschenschaftern an den Universitäten).
        Ebenso sind Kulturveranstaltungen erwünscht (z.B. Lesungen, Theater,
        Filmvorführungen, Campusrundgang, Würdigung der Gedenktafel für die
        Weiße Rose im Audimax, Ausstellung zur Feuerzangenbowle). 
    \end{enumerate}

\end{document}
